\documentclass{beamer}

\usepackage[utf8]{inputenc}
\input{../helper}

\title{Introdução à Computação \\ compilando primeiro programa em C}
\author{Alexandre Rademaker}
\date{}

\begin{document}

\begin{frame}
 \maketitle
\end{frame}

\begin{frame}[fragile,allowframebreaks]{Primeiro programa em C}

  \lstinputlisting[caption=src/hello0.c,style=normalc]{src/hello0.c}

  \framebreak

  Comentários servem para documentar o código.

  A importação do \ilcode{stdio.h}. Um arquivo de cabeçalho (header)
  tem extensão \ilcode{.h} e contém declarações de funções e
  definições de macros compartilhadas entre vários arquivos
  \ilcode{.c}.

  A função main é o `entry point' em um programa em C. Ela pode
  receber parâmetros que permitem a comunicação entre o programa em C
  e o sistema operacional.

  Funções recebem lista de argumentos entre parênteses, e retornam
  valores ou apenas produzem `efeitos colaterais'.

  \framebreak

  O nome
  \href{https://man7.org/linux/man-pages/man3/printf.3.html}{printf} é
  de `formatted string'.

  O ponto-e-vírgula no final de todos os comandos.

  O comando \ilcode{return} recebe um valor inteiro que será retornado
  ao SO no final da execução, zero é `sem erro'.

  A identação organiza o código em blocos lógicos e destacando a
  estrutura de controle do programa e tornando-o mais legível.

  Vamos editar no VS Code, Replit ou GitHub Codespaces e compilar!

  \framebreak  

  \includegraphics[width=\textwidth]{vs-code.png}

\end{frame}


\begin{frame}[fragile,allowframebreaks]{Compilação}

  O compilador \ilcode{clang} traduz código C em um executável. Na
  forma mais simples:

  \begin{lstlisting}[style=normalbash]
    > clang hello0.c
    > ./a.out
  \end{lstlisting}

  Funciona, mas sem nenhuma verificação adicional. O compilador aceita
  código que pode conter problemas. Para maior controle, passamos
  flags:

  \begin{lstlisting}[style=normalbash]
    > clang -std=c17 -Wall -Wextra -Wpedantic -Werror hello0.c -o hello
    > ./hello
  \end{lstlisting}

  Em alguns ambientes, \ilcode{cc} ou \ilcode{gcc} ao invés de
  \ilcode{clang}.

  \framebreak

  As flags do \ilcode{clang} nos dão controle sobre a compilação:

  \begin{description}
  \item[\texttt{-std=c17}] versão do padrão C a ser usada
  \item[\texttt{-Wall -Wextra -Wpedantic}] ativam avisos (warnings)
    que ajudam a encontrar erros no código
  \item[\texttt{-Werror}] trata avisos como erros, forçando código
    mais rigoroso
  \item[\texttt{-o hello}] define o nome do executável de saída (sem
    esta flag, o padrão é \ilcode{a.out})
  \end{description}

  \framebreak

  \begin{lstlisting}[style=code]
    > clang src/cents1.c
    ...
    clang: error: linker command failed with exit code 1 (...)

    > clang src/cents1.c -lm
  \end{lstlisting}

  Em alguns casos poderão ser necessários parâmetros adicionais para o
  compilador.

  O programa usa a função \ilcode{round} da biblioteca matemática
  (\ilcode{math.h}). No Linux, essa biblioteca (\ilcode{libm}) é
  separada e precisa ser vinculada explicitamente com \ilcode{-lm}. No
  macOS, a \ilcode{libm} faz parte da biblioteca do sistema
  (\ilcode{libSystem}) e é vinculada automaticamente.

  \framebreak

  O comando
  \href{https://www.gnu.org/software/make/}{\texttt{make}} é um
  utilitário que automatiza a compilação. Ele executa uma sequência de
  comandos (tipicamente chamadas ao compilador) seguindo regras
  definidas em um arquivo chamado \ilcode{Makefile}.

  \begin{lstlisting}[style=normalbash]
    > make hello0
    > ./hello0
  \end{lstlisting}

  Quando \textbf{não existe} um \ilcode{Makefile}, o \ilcode{make} tem
  um comportamento \emph{default}: para \ilcode{make hello0}, ele
  procura \ilcode{hello0.c} e chama o compilador com regras implícitas
  (ex: \ilcode{cc hello0.c -o hello0}).

  A limitação é justamente a falta de controle: com o \ilcode{make}
  default, não escolhemos as flags passadas ao compilador.

  \framebreak

  Compilar o código na verdade significa:

  \begin{itemize}
  \item pré-processamento (pre-processing)
  \item compilação (compiling)
  \item montagem (assembling)
  \item vinculação (linking)   
  \end{itemize}

  \framebreak
  
  \lstinputlisting[style=normalc,caption=src/cents1.c]{src/cents1.c}

  \framebreak

  \begin{lstlisting}[style=normalc]
    ...
    double round(double n);
    ...
    int scanf(string format, ...);
    int printf(string format, ...);
    ...

    int main(void) {
      float amount;
      printf("Dollar amount: ");
      scanf("%f", &amount);

      int cents = round(amount * 100);
      printf("Cents: %i\n", cents);
    }
  \end{lstlisting}
  
  \framebreak
  
  O código C será transformado em código do processador (arquitetura)
  específica sendo usada.\vspace{.5cm}

  \begin{lstlisting}[style=normalbash]
    > clang -S -mllvm --x86-asm-syntax=intel hello0.c
  \end{lstlisting}

  Vide \ilcode{src/hello0.s}

  \framebreak

  \begin{lstlisting}
    hello0.a  libm.a      libc.a
    
    01111111  0100010101  001100010001
    00000010  0000000100  000001000000
    00000000  0000000000  000000000000
    00000000  0000000000  000000000000
    00000001  0000000000  111110000000
    00000001  0000000000  000100000000
    00000000  0000000000  000000001000
    00000000  0000000000  000000000000
    00000000  0000000000  000000010000
    00000000  0000000000  000000000000
    10100000  0000001000  000010001000
    00000000  0000000000  000000000000
  \end{lstlisting}

  \framebreak

  \begin{lstlisting}
    hello0
    
    01111111010001010100110001000110
    00000010000000010000000100000000
    00000000000000000000000000000000
    00000000000000000000000000000000
    00000001000000000011111000000000
    00000001000000000000000000000000
    00000000000000000000000000000000
    00000000000000000000000000000000
    00000000000000000000000000000000
    00000000000000000000000000000000
    10100000000000100000000000000000
    00000000000000000000000000000000
  \end{lstlisting}

\end{frame}

\begin{frame}[fragile,allowframebreaks]{Erros}

  Erros ocorrem, e precisamos aprender a lidar com eles!

  \begin{itemize}
  \item O que ocorre se removermos o \ilcode{\\n}?
  \item O que ocorre se quebrarmos a string?
  \item O que ocorre se esquecermos o ponto-e-vírgula?
  \item O que ocorre se removermos a biblioteca?
  \end{itemize}
  
  O símbolo \ilcode{\\n} e o `escape'.

\end{frame}


\begin{frame}[allowframebreaks]{Um bom código C}

  \begin{description}
  \item[correção] seu código resolve o problema de forma correta?

  \item[design] quão eficiente é seu código ou seu \emph{algorítmo}?

  \item[estilo] quão bem escrito, formatado e legível é seu código?
  \end{description}

  \framebreak
  
  \begin{block}{}
    ``Always code as if the guy who ends up maintaining your code will
    be a violent psychopath who knows where you live'' -- John Woods
  \end{block}
  
  \begin{block}{}
    ``Programs must be written for people to read, and only
    incidentally for machines to execute.''  -- Harold Abelson,
    Structure and Interpretation of Computer Programs
  \end{block}

  Mantenha seu código legível! Mais
  \href{https://www.goodreads.com/quotes/tag/programming}{quotes}!

\end{frame}  



\end{document}

%%% Local Variables:
%%% mode: latex
%%% TeX-master: t
%%% End:
